\section{Discussion and conclusion}\label{sec:short-discussion}
The firm-stage and political mechanisms should not be conflated. Supporting an additional standard has one-time cost $F$ and may leave residual marginal cost $d$; where outsiders adopt unilaterally, consumers in member markets can benefit while protected domestic producers lose profits. Governments count their own consumers and producers, not the foreign adopter. The residual rent $R(d,v)$ therefore changes incentives for \emph{formal} harmonization without changing the formal partition itself. Conditional on a fixed outsider-only continuation and the canonical market structure, lowering $d$ reduces the member's relative preference for $SU$; it is not an unconditional prediction across technologies or countries.

Our results do not imply that international standardization is the globally first-best policy, nor that the positive-$d$ political ranking determines a stable coalition structure. The full-bypass welfare and blocking comparisons concern only the five formal partitions under specified rules. The distinction between technology and policy requires care. A regulator may observe an increase in products that technically comply with multiple standards and conclude that regional arrangements are becoming redundant; the model shows that the national-welfare incentive to negotiate a broader agreement also depends on the displaced domestic operating profits and on disadvantages that remain in outside destinations. As a result, observing private standard adoption alone cannot identify the sign of the government preference effect. In the canonical ranking-reversal region, this is a strictly national-welfare rather than a global-planner comparison. The model uses segmented linear demand, one firm per country, a singleton-zero network benefit and binary adoption; the identified preference reversal is not portable even to the pre-specified differentiated-demand environment. Empirical use would require separate measures of adoption setup costs, remaining per-unit compliance costs, demand-side network benefits, home producer rents, and whether interoperability is reciprocal. No empirically identified policy intervention is claimed. This conditionality is itself the economic lesson: private adaptation can diminish technological incompatibility without making broader public harmonization either automatically more or automatically less attractive.
